\section*{Capítulo \ref{ch:g7:sediments-and-fossils} --- Sedimentos y fósiles}

\begin{solution}{exo:g7:sediments-and-fossils:1}
(a) \omterm{def:g7:sediments-and-fossils:fossil}{molde}; (b) \omterm{def:g7:sediments-and-fossils:fossil}{icnofósil}; (c) \omterm{def:g7:sediments-and-fossils:sediment}{sedimento}; (d) \omterm{def:g7:sediments-and-fossils:fossil}{contramolde}; (e) \omterm{def:g7:sediments-and-fossils:layer}{estrato};
(f) \omterm{def:g7:sediments-and-fossils:fossil}{fósil}.
\end{solution}

\begin{solution}{exo:g7:sediments-and-fossils:2}
(b) el animal muere; (d) el \omterm{def:g3:stones-pebbles-sand:sand}{barro} cubre la concha deprisa; (c) se acumulan
capas nuevas y la concha se convierte en piedra; (a) la \omterm{def:g5:water-shapes-the-land:erosion}{erosión} desgasta
las capas de encima.
\end{solution}

\begin{solution}{exo:g7:sediments-and-fossils:3}
(a) Un \omterm{def:g2:water-on-earth:ocean}{mar} cálido, poco profundo y transparente. (b) Aguas poco profundas,
en la orilla de un \omterm{def:g1:landscapes-around-us:water}{lago}, de un \omterm{def:g1:landscapes-around-us:water}{río} o de un \omterm{def:g2:water-on-earth:ocean}{mar}. (c) Un bosque pantanoso o
una marisma. (d) Tierra firme, con plantas (y libélulas volando), cerca de
un agua donde se depositaba el \omterm{def:g3:stones-pebbles-sand:sand}{barro}.
\end{solution}

\begin{solution}{exo:g7:sediments-and-fossils:4}
(a) Falso: la mayoría de los seres vivos no dejan ningún \omterm{def:g7:sediments-and-fossils:fossil}{fósil}; se pudren
o se disuelven antes de quedar enterrados. (b) Verdadero: es un \omterm{def:g7:sediments-and-fossils:fossil}{icnofósil}.
(c) Falso: los granos más grandes y pesados se depositan primero, y el
\omterm{def:g3:stones-pebbles-sand:sand}{barro} fino, al final. (d) Verdadero.
\end{solution}

\begin{solution}{exo:g7:sediments-and-fossils:5}
(a) Arenisca \qty{120}{m}, caliza \qty{180}{m}, lutita \qty{100}{m}.
(b) $120 + 180 + 100 = \qty{400}{m}$. (c) $180 \div 400 = 0.45$:
\textbf{45~\%}. (d) La \omterm{def:g5:water-shapes-the-land:erosion}{erosión} ha desgastado su techo: lo que queda no es
la capa entera, y lo que queda cambia de un lugar a otro.
\end{solution}

\begin{solution}{exo:g7:sediments-and-fossils:6}
(a) De abajo arriba: \omterm{def:g3:stones-pebbles-sand:pebble}{grava}, después \omterm{def:g3:stones-pebbles-sand:sand}{arena} y, encima, \omterm{def:g3:stones-pebbles-sand:sand}{barro}, bajo el agua
transparente. (b) Los granos diminutos de \omterm{def:g3:stones-pebbles-sand:sand}{barro} se hunden muy despacio:
siguen cayendo al cabo de una hora. (c) Una segunda serie de capas encima
de la primera: otra vez \omterm{def:g3:stones-pebbles-sand:sand}{arena} y luego \omterm{def:g3:stones-pebbles-sand:sand}{barro}. (d) La \omterm{def:g3:stones-pebbles-sand:pebble}{grava}, en el fondo del
todo.
\end{solution}

\begin{solution}{exo:g7:sediments-and-fossils:7}
Más posibilidades: (b), una concha dura, enterrada deprisa bajo \omterm{def:g3:stones-pebbles-sand:sand}{barro} fino
en agua tranquila. Después, (d): la concha es dura y queda enterrada, pero
una \omterm{def:g5:water-shapes-the-land:flood}{crecida} también puede moverla y romperla. Menos posibilidades: (a), un
cuerpo blando sin partes duras, sobre \omterm{def:g3:stones-pebbles-sand:rock}{rocas} donde no se deposita ningún
\omterm{def:g3:stones-pebbles-sand:sand}{barro} que lo cubra. (c) no está mucho mejor: el hueso se queda al
descubierto y se descompone y se desgasta.
\end{solution}

\begin{solution}{exo:g7:sediments-and-fossils:8}
(a) Un hueco con la forma de la concha impresa en sus paredes: un \omterm{def:g7:sediments-and-fossils:fossil}{molde}.
(b) Un \omterm{def:g7:sediments-and-fossils:fossil}{contramolde}, una copia de piedra de la concha. (c) No: la concha se
ha disuelto; el \omterm{def:g7:sediments-and-fossils:fossil}{contramolde} está hecho de \omterm{def:g6:rocks-and-minerals:mineral}{minerales} nuevos. (d) El agua
pasa fácilmente entre los granos gruesos de la \omterm{def:g3:stones-pebbles-sand:sand}{arena}, así que puede llegar
a la concha y disolverla; en el \omterm{def:g3:stones-pebbles-sand:sand}{barro} fino pasa con mucha más dificultad.
\end{solution}

\begin{solution}{exo:g7:sediments-and-fossils:9}
(a) $1.6 \times 1000 = \qty{1600}{m}$. (b) $1600 \div 80 = 20$ capas.
(c) $160 \div 1600 = 0.1$: \textbf{10~\%}.
\end{solution}

\begin{solution}{exo:g7:sediments-and-fossils:10}
(a) No: un pie solo deja huella en un suelo blando, en \omterm{def:g3:stones-pebbles-sand:sand}{arena} mojada o en
\omterm{def:g3:stones-pebbles-sand:sand}{barro}; el suelo se convirtió en \omterm{def:g3:stones-pebbles-sand:rock}{roca} más tarde. (b) La ola o la \omterm{def:g5:water-shapes-the-land:flood}{crecida}
siguiente las borra. (c) Un \omterm{def:g7:sediments-and-fossils:sediment}{sedimento} nuevo tuvo que cubrirlas y
rellenarlas deprisa, antes de que se borraran. (d) Un \omterm{def:g7:sediments-and-fossils:fossil}{icnofósil}: muestra lo
que hizo el animal, no una parte de su cuerpo.
\end{solution}

\begin{solution}{exo:g7:sediments-and-fossils:11}
(a) $20 + 5 = \qty{25}{mm}$. (b) $1000 \div 25 = 40$ pares: el testigo
registra \textbf{40 \omterm{def:g5:water-shapes-the-land:flood}{crecidas}}. (c) La \omterm{def:g3:stones-pebbles-sand:sand}{arena} es pesada: se deposita poco
después de que el \omterm{def:g1:landscapes-around-us:water}{río} entre en el \omterm{def:g1:landscapes-around-us:water}{lago}, cerca de la \omterm{def:g5:water-shapes-the-land:source-mouth}{desembocadura}; solo
una pequeña parte llega lejos. (d) Lejos, en el agua ya solo queda \omterm{def:g3:stones-pebbles-sand:sand}{barro}
fino: se deposita allí despacio, cuando el agua está tranquila.
\end{solution}

\begin{solution}{exo:g7:sediments-and-fossils:12}
(a) Las conchas son duras: duran lo suficiente para quedar enterradas, y
muchos animales con concha vivían en el fondo del \omterm{def:g2:water-on-earth:ocean}{mar}, donde el \omterm{def:g3:stones-pebbles-sand:sand}{barro} se
deposita sobre ellos. (b) Los cuerpos blandos se pudren antes de quedar
enterrados. (c) No: los animales blandos casi nunca dejan \omterm{def:g7:sediments-and-fossils:fossil}{fósiles}, así que
la lista que dan los \omterm{def:g7:sediments-and-fossils:fossil}{fósiles} está incompleta; muestra sobre todo los
animales con concha o con huesos.
\end{solution}

\begin{solution}{pb:g7:sediments-and-fossils:1}
\textbf{1.} Arenisca, caliza, lutita, creta (la primera, en la base de la
columna).

\textbf{2.} \qty{120}{m}, \qty{180}{m} y \qty{100}{m}; en total,
\qty{400}{m}.

\textbf{3.} $180 \div 400 = 0.45$: \textbf{45~\%}.

\textbf{4.} Su techo se ha desgastado: la parte superior de la capa ha
desaparecido, y nadie puede medir lo que ya no está.

\textbf{5.} En aguas poco profundas, en la orilla de un \omterm{def:g1:landscapes-around-us:water}{lago}, de un \omterm{def:g1:landscapes-around-us:water}{río} o
de un \omterm{def:g2:water-on-earth:ocean}{mar}: las rizaduras se forman sobre la \omterm{def:g3:stones-pebbles-sand:sand}{arena} en aguas poco profundas,
y las huellas se quedan en \omterm{def:g3:stones-pebbles-sand:sand}{arena} o \omterm{def:g3:stones-pebbles-sand:sand}{barro} blandos y mojados, a menudo en
una \omterm{def:g1:landscapes-around-us:water}{costa}.

\textbf{6.} En un \omterm{def:g2:water-on-earth:ocean}{mar} cálido, poco profundo y transparente: los corales.

\textbf{7.} Tranquila: el \omterm{def:g3:stones-pebbles-sand:sand}{barro} fino solo se deposita donde el agua está casi quieta.

\textbf{8.} Al principio, el lugar era una \omterm{def:g1:landscapes-around-us:water}{costa}, con aguas poco profundas
y \omterm{def:g3:stones-pebbles-sand:sand}{arena} blanda por la que caminaban animales. Después, el \omterm{def:g2:water-on-earth:ocean}{mar} lo cubrió:
un \omterm{def:g2:water-on-earth:ocean}{mar} cálido y poco profundo donde crecían corales y se acumulaban
conchas. Más tarde, el agua se volvió más tranquila, quizá más profunda, y
solo se depositó \omterm{def:g3:stones-pebbles-sand:sand}{barro} fino, con las conchas de los amonites, animales
emparentados con los calamares y los pulpos.

\textbf{9.} $180 \times 100 = 18\,000$~cm.

\textbf{10.} $120~\text{m} = 12\,000$~cm, y $12\,000 \times 20 =
240\,000$~años.

\textbf{11.} $1~\text{m} = 100$~cm, y $100 \times 100 = 10\,000$~años.

\textbf{12.} $180 \times 10\,000 = 1\,800\,000$: unos
\textbf{1\,800\,000 años}, casi dos millones de años. Es
$1\,800\,000 \div 240\,000 = 7.5$ veces más que la arenisca.
\end{solution}
